\section{Step 9: Advanced Features}

Argument parsing already collected values for custom headers, POST bodies,
output files, basic auth, and HEAD requests. This step \emph{uses} those
values when building and sending the HTTP message, and when writing the
response.

\subsection{Method Selection}

\begin{lstlisting}[style=snippet,language=C]
const char *method = head_request ? "HEAD"
                     : (post_data[0] ? "POST" : "GET");
\end{lstlisting}

Priority: explicit \texttt{-I}/\texttt{--head} wins; otherwise any
\texttt{-d} data implies \texttt{POST}; default remains \texttt{GET}.

\subsection{Request Construction}

Optional header blocks are injected as pre-formatted strings ending in
\texttt{\textbackslash r\textbackslash n} (see the argument parser):

\begin{lstlisting}[style=snippet,language=C]
int len = snprintf(req, sizeof(req),
    "%s %s HTTP/1.1\r\n"
    "Host: %s\r\n"
    "User-Agent: %s\r\n"
    "%s"   /* extra_headers */
    "%s"   /* auth_header   */
    "Accept: */*\r\n"
    "Connection: close\r\n",
    method, u.path, u.host, user_agent,
    extra_headers, auth_header);

if (post_data[0]) {
    len += snprintf(req + len, sizeof(req) - len,
        "Content-Length: %zu\r\n\r\n%s",
        strlen(post_data), post_data);
} else {
    strcat(req + len, "\r\n");
}
\end{lstlisting}

\subsection{Output File (\texttt{-o})}

\begin{lstlisting}[style=snippet,language=C]
FILE *out = stdout;
if (output_file[0]) {
    out = fopen(output_file, "w");
    if (!out) { perror("fopen"); return EXIT_FAILURE; }
}
/* ... fwrite(buf, 1, n, out); ... */
if (out != stdout) fclose(out);
\end{lstlisting}

\subsection{Basic Authentication (\texttt{-u}) with Base64}

HTTP Basic auth expects
\texttt{Authorization: Basic <base64(user:pass)>}.
We accept curl-style \texttt{USER:PASS} and encode it ourselves with a small
RFC~4648 helper---no external process, no extra libraries.

\begin{lstlisting}[style=snippet,language=C]
static inline int base64_encode(const void *data, size_t len,
                                char *out, size_t out_sz)
{
    static const char T[] =
        "ABCDEFGHIJKLMNOPQRSTUVWXYZabcdefghijklmnopqrstuvwxyz"
        "0123456789+/";
    const unsigned char *s = data;
    size_t o = 0;

    for (size_t i = 0; i < len; i += 3) {
        unsigned v = (unsigned)s[i] << 16;
        int n = 1;
        if (i + 1 < len) { v |= (unsigned)s[i + 1] << 8; n = 2; }
        if (i + 2 < len) { v |= (unsigned)s[i + 2];      n = 3; }
        if (o + 4 >= out_sz) return -1;
        out[o++] = T[(v >> 18) & 63];
        out[o++] = T[(v >> 12) & 63];
        out[o++] = (n >= 2) ? T[(v >> 6) & 63] : '=';
        out[o++] = (n >= 3) ? T[v & 63]         : '=';
    }
    if (o >= out_sz) return -1;
    out[o] = '\0';
    return (int)o;
}
\end{lstlisting}

When parsing \texttt{-u} / \texttt{--user}:

\begin{lstlisting}[style=snippet,language=C]
char b64[200];
const char *cred = argv[++i];
if (base64_encode(cred, strlen(cred), b64, sizeof(b64)) < 0) {
    fprintf(stderr, "credentials too long\n");
    return EXIT_FAILURE;
}
snprintf(auth_header, sizeof(auth_header),
         "Authorization: Basic %s\r\n", b64);
\end{lstlisting}

The encoder walks the input three bytes at a time, emits four alphabet
characters, and writes \texttt{=} padding when the final group is short.
That is enough for credential strings and keeps the whole client
dependency-free beyond OpenSSL (which we already link for HTTPS).

\FullSourceStart{urlgetstep09}
\begin{lstlisting}[style=oldcode,language=C,firstnumber=1]
/* urlget.c -- Step 9: advanced request features */
#include <stdio.h>
#include <stdlib.h>
#include <string.h>
#include <stdbool.h>
#include <stdint.h>
#include <unistd.h>
#include <sys/types.h>
#include <sys/socket.h>
#include <netdb.h>
#include <openssl/ssl.h>
#include <openssl/err.h>
\end{lstlisting}
\begin{lstlisting}[style=newcode,language=C,firstnumber=last]
/* RFC 4648 Base64 for HTTP Basic auth (-u USER:PASS). */
static inline int base64_encode(const void *data, size_t len,
                                char *out, size_t out_sz)
{
    static const char T[] =
        "ABCDEFGHIJKLMNOPQRSTUVWXYZabcdefghijklmnopqrstuvwxyz"
        "0123456789+/";
    const unsigned char *s = data;
    size_t o = 0;

    for (size_t i = 0; i < len; i += 3) {
        unsigned v = (unsigned)s[i] << 16;
        int n = 1;
        if (i + 1 < len) { v |= (unsigned)s[i + 1] << 8; n = 2; }
        if (i + 2 < len) { v |= (unsigned)s[i + 2];      n = 3; }
        if (o + 4 >= out_sz) return -1;
        out[o++] = T[(v >> 18) & 63];
        out[o++] = T[(v >> 12) & 63];
        out[o++] = (n >= 2) ? T[(v >> 6) & 63] : '=';
        out[o++] = (n >= 3) ? T[v & 63]         : '=';
    }
    if (o >= out_sz) return -1;
    out[o] = '\0';
    return (int)o;
}
\end{lstlisting}
\begin{lstlisting}[style=oldcode,language=C,firstnumber=last]
typedef struct {
    bool https;
    char host[256];
    uint16_t port;
    char path[512];
} url_t;

static inline int parse_url(const char *arg, url_t *u)
{
    char buf[1024];
    if (strlen(arg) >= sizeof(buf)) return -1;
    strcpy(buf, arg);
    char *p = buf;
    u->https = false;
    if (strncmp(p, "https://", 8) == 0) { u->https = true; p += 8; }
    else if (strncmp(p, "http://", 7) == 0) { p += 7; }
    char *sep = strpbrk(p, "/?");
    char *hostport = p;
    if (sep) {
        if (*sep == '?') {
            size_t ql = strlen(sep);
            if (ql + 2 > sizeof(u->path)) return -1;
            u->path[0] = '/';
            memcpy(u->path + 1, sep, ql + 1);
        } else {
            if (strlen(sep) >= sizeof(u->path)) return -1;
            strcpy(u->path, sep);
        }
        *sep = '\0';
    } else strcpy(u->path, "/");
    char *host = u->host;
    size_t hmax = sizeof(u->host) - 1;
    uint16_t def = u->https ? 443 : 80;
    u->port = def;
    if (*hostport == '[') {
        char *cl = strchr(hostport, ']');
        if (!cl) return -1;
        size_t hl = cl - (hostport + 1);
        if (hl >= hmax) return -1;
        memcpy(host, hostport + 1, hl);
        host[hl] = '\0';
        if (*(cl + 1) == ':')
            u->port = (uint16_t)strtol(cl + 2, NULL, 10);
    } else {
        char *col = strrchr(hostport, ':');
        if (col) {
            size_t hl = col - hostport;
            if (hl >= hmax) return -1;
            memcpy(host, hostport, hl);
            host[hl] = '\0';
            u->port = (uint16_t)strtol(col + 1, NULL, 10);
        } else {
            if (strlen(hostport) >= hmax) return -1;
            strcpy(host, hostport);
        }
    }
    if (u->port == 0) u->port = def;
    return 0;
}

static int socks5_connect(int sock, const char *host, uint16_t port)
{
    unsigned char buf[512];
    size_t hlen = strlen(host);
    if (hlen > 255) return -1;
    buf[0] = 0x05; buf[1] = 0x01; buf[2] = 0x00;
    if (write(sock, buf, 3) != 3) return -1;
    if (read(sock, buf, 2) != 2 || buf[0] != 0x05 || buf[1] != 0x00)
        return -1;
    buf[0] = 0x05; buf[1] = 0x01; buf[2] = 0x00; buf[3] = 0x03;
    buf[4] = (unsigned char)hlen;
    memcpy(buf + 5, host, hlen);
    buf[5 + hlen] = port >> 8;
    buf[6 + hlen] = port & 0xff;
    if (write(sock, buf, 7 + hlen) != (ssize_t)(7 + hlen)) return -1;
    if (read(sock, buf, 4) != 4 || buf[1] != 0x00) return -1;
    return 0;
}

static inline bool https_setup(int sock, const char *hostname,
                               SSL_CTX **out_ctx, SSL **out_ssl)
{
    SSL_CTX *ctx = SSL_CTX_new(TLS_client_method());
    if (!ctx) { ERR_print_errors_fp(stderr); return false; }
    SSL_CTX_set_verify(ctx, SSL_VERIFY_NONE, NULL);
    SSL *ssl = SSL_new(ctx);
    if (!ssl) { SSL_CTX_free(ctx); return false; }
    SSL_set_fd(ssl, sock);
    SSL_set_tlsext_host_name(ssl, hostname);
    if (SSL_connect(ssl) <= 0) {
        ERR_print_errors_fp(stderr);
        SSL_free(ssl); SSL_CTX_free(ctx);
        return false;
    }
    *out_ctx = ctx;
    *out_ssl = ssl;
    return true;
}

static inline void usage(const char *progname)
{
    fprintf(stderr,
        "Usage: %s [options] <URL>\n\n"
        "Options:\n"
        "  -A, --user-agent STR     Set User-Agent\n"
        "  -T, --tor [proxy]        Use Tor (default 127.0.0.1:9050)\n"
        "  -H, --header STR         Add custom header (repeatable)\n"
        "  -d, --data DATA          POST data (@file supported)\n"
        "  -o, --output FILE        Write output to file\n"
        "  -u, --user USER:PASS     Basic authentication\n"
        "  -I, --head               HEAD request\n"
        "      --no-location        Do not follow redirects\n"
        "      --max-redirects N    Max redirects (default: 10)\n",
        progname);
    exit(EXIT_FAILURE);
}

int main(int argc, char *argv[])
{
    const char *user_agent =
        "Mozilla/5.0 (X11; Linux x86_64; rv:128.0) "
        "Gecko/20100101 Firefox/128.0";
    const char *default_tor = "127.0.0.1:9050";
    const char *url_arg = NULL;
    const char *tor_proxy = NULL;
    bool use_tor = false;
    bool follow_redirects = true;
    int max_redirects = 10;
    bool head_request = false;
    char extra_headers[4096] = {0};
    char post_data[8192] = {0};
    char output_file[256] = {0};
    char auth_header[256] = {0};

    for (int i = 1; i < argc; i++) {
        if (strcmp(argv[i], "-A") == 0 ||
            strcmp(argv[i], "--user-agent") == 0) {
            if (i + 1 < argc) user_agent = argv[++i];
            else usage(argv[0]);
        }
        else if (strcmp(argv[i], "-T") == 0 ||
                 strcmp(argv[i], "--tor") == 0) {
            use_tor = true;
            if (i + 1 < argc && argv[i + 1][0] != '-' &&
                strchr(argv[i + 1], ':')) {
                tor_proxy = argv[i + 1];
                i++;
            }
        }
        else if (strcmp(argv[i], "-H") == 0 ||
                 strcmp(argv[i], "--header") == 0) {
            if (i + 1 < argc) {
                strncat(extra_headers, argv[++i],
                        sizeof(extra_headers) -
                        strlen(extra_headers) - 3);
                strcat(extra_headers, "\r\n");
            } else usage(argv[0]);
        }
        else if (strcmp(argv[i], "-d") == 0 ||
                 strcmp(argv[i], "--data") == 0) {
            if (i + 1 < argc) {
                const char *val = argv[++i];
                if (val[0] == '@') {
                    FILE *f = fopen(val + 1, "r");
                    if (f) {
                        fread(post_data, 1,
                              sizeof(post_data) - 1, f);
                        fclose(f);
                    }
                } else
                    strncpy(post_data, val, sizeof(post_data) - 1);
            } else usage(argv[0]);
        }
        else if (strcmp(argv[i], "-o") == 0 ||
                 strcmp(argv[i], "--output") == 0) {
            if (i + 1 < argc)
                strncpy(output_file, argv[++i],
                        sizeof(output_file) - 1);
            else usage(argv[0]);
        }
\end{lstlisting}
\begin{lstlisting}[style=newcode,language=C,firstnumber=last]
        else if (strcmp(argv[i], "-u") == 0 ||
                 strcmp(argv[i], "--user") == 0) {
            if (i + 1 < argc) {
                char b64[200];
                const char *cred = argv[++i];
                if (base64_encode(cred, strlen(cred), b64,
                                  sizeof(b64)) < 0) {
                    fprintf(stderr, "credentials too long\n");
                    return EXIT_FAILURE;
                }
                snprintf(auth_header, sizeof(auth_header),
                         "Authorization: Basic %s\r\n", b64);
            } else usage(argv[0]);
        }
\end{lstlisting}
\begin{lstlisting}[style=oldcode,language=C,firstnumber=last]
        else if (strcmp(argv[i], "-I") == 0 ||
                 strcmp(argv[i], "--head") == 0)
            head_request = true;
        else if (strcmp(argv[i], "--no-location") == 0)
            follow_redirects = false;
        else if (strcmp(argv[i], "--max-redirects") == 0 &&
                 i + 1 < argc)
            max_redirects = atoi(argv[++i]);
        else if (url_arg == NULL)
            url_arg = argv[i];
    }

    if (url_arg == NULL) usage(argv[0]);
    if (use_tor && tor_proxy == NULL) tor_proxy = default_tor;
\end{lstlisting}
\begin{lstlisting}[style=newcode,language=C,firstnumber=last]
    FILE *out = stdout;
    if (output_file[0]) {
        out = fopen(output_file, "w");
        if (!out) { perror("fopen"); return EXIT_FAILURE; }
    }

    int redirect_count = 0;
    char current_url[1024];
    strncpy(current_url, url_arg, sizeof(current_url) - 1);

    while (redirect_count <= max_redirects) {
        url_t u = {0};
        if (parse_url(current_url, &u) != 0) {
            fprintf(stderr, "Failed to parse URL\n");
            return EXIT_FAILURE;
        }

        const char *connect_host = use_tor ? tor_proxy : u.host;
        uint16_t connect_port = use_tor ? 9050 : u.port;
        char pbuf[256];
        if (use_tor) {
            strncpy(pbuf, tor_proxy, sizeof(pbuf) - 1);
            char *colon = strrchr(pbuf, ':');
            if (colon) {
                *colon = '\0';
                connect_port =
                    (uint16_t)strtol(colon + 1, NULL, 10);
                connect_host = pbuf;
            }
        }

        struct addrinfo hints = {0};
        hints.ai_family   = AF_UNSPEC;
        hints.ai_socktype = SOCK_STREAM;
        char ps[6];
        snprintf(ps, sizeof(ps), "%u", connect_port);

        struct addrinfo *res = NULL;
        if (getaddrinfo(connect_host, ps, &hints, &res) != 0) {
            perror("getaddrinfo");
            return EXIT_FAILURE;
        }

        int s = -1;
        for (struct addrinfo *a = res; a; a = a->ai_next) {
            s = socket(a->ai_family, a->ai_socktype, a->ai_protocol);
            if (s < 0) continue;
            if (connect(s, a->ai_addr, a->ai_addrlen) == 0) break;
            close(s); s = -1;
        }
        freeaddrinfo(res);
        if (s < 0) { perror("connect"); return EXIT_FAILURE; }

        if (use_tor && socks5_connect(s, u.host, u.port) != 0) {
            close(s);
            return EXIT_FAILURE;
        }

        SSL_CTX *ctx = NULL;
        SSL *ssl = NULL;
        if (u.https && !https_setup(s, u.host, &ctx, &ssl)) {
            close(s);
            return EXIT_FAILURE;
        }

        const char *method = head_request ? "HEAD"
                             : (post_data[0] ? "POST" : "GET");

        char req[4096];
        int len = snprintf(req, sizeof(req),
            "%s %s HTTP/1.1\r\n"
            "Host: %s\r\n"
            "User-Agent: %s\r\n"
            "%s"
            "%s"
            "Accept: */*\r\n"
            "Connection: close\r\n",
            method, u.path, u.host, user_agent,
            extra_headers, auth_header);

        if (post_data[0]) {
            len += snprintf(req + len, sizeof(req) - len,
                "Content-Length: %zu\r\n\r\n%s",
                strlen(post_data), post_data);
        } else {
            strcat(req + len, "\r\n");
        }

        if (u.https)
            SSL_write(ssl, req, strlen(req));
        else
            write(s, req, strlen(req));

        char buf[8192];
        ssize_t n;
        while ((n = u.https
                    ? SSL_read(ssl, buf, sizeof(buf))
                    : read(s, buf, sizeof(buf))) > 0)
            fwrite(buf, 1, n, out);

        if (u.https) {
            if (ssl) { SSL_shutdown(ssl); SSL_free(ssl); }
            if (ctx) SSL_CTX_free(ctx);
        }
        close(s);

        if (follow_redirects && strstr(buf, "Location:") &&
            redirect_count < max_redirects) {
            char location[1024] = {0};
            char *loc = strstr(buf, "Location:");
            if (loc) sscanf(loc, "Location: %1023s", location);
            if (location[0]) {
                redirect_count++;
                strncpy(current_url, location,
                        sizeof(current_url) - 1);
                continue;
            }
        }
        break;
    }

    if (out != stdout) fclose(out);
\end{lstlisting}
\begin{lstlisting}[style=oldcode,language=C,firstnumber=last]
    return 0;
}
\end{lstlisting}
\FullSourceStop

\begin{lstlisting}[style=snippet,language=bash]
./urlget -I https://example.com/
./urlget -H "Accept-Language: en" https://example.com/
./urlget -d "name=Ada" -o resp.txt https://httpbin.org/post
# curl-style USER:PASS --- encoded inside urlget
./urlget -u 'user:pass' https://httpbin.org/basic-auth/user/pass
\end{lstlisting}
